\section{Numerical Validation and Discussion}
\label{sec:experiments}

The experiments examine whether LLM guidance improves charging optimization
at a fixed simulation budget and which part of the guidance accounts for the
gain. Comparisons with other optimizers are followed by initialization and
component studies. We also compare LLM backends and inspect representative
charging protocols from the computed Pareto set.

\subsection{Experimental Setup}

The main case uses the Chen2020 parameterization of the LG INR21700-M50
cell~\cite{ref8}. A second case uses Ecker2015~\cite{ref9} to examine
performance with a different battery parameterization. The objectives are
charging time $t_c$, peak temperature rise $T_p$, and aging indicator $Q_s$.
Tables~\ref{tab:battery_settings} and~\ref{tab:algorithm_settings} summarize
the battery and algorithm settings.

The main comparison allocates 56 battery-model evaluations to each method
and uses five matched trials. For the BO methods, this budget consists of six
initial evaluations and 50 subsequent iterations. Hypervolume (HV) measures
the quality of the feasible nondominated set; larger values are better under
the same objective definitions and reference point. All methods use the same
HV calculation within each battery parameterization. Final HV is reported as
mean $\pm$ sample standard deviation.

\begin{table}[!t]
\caption{Battery and Simulator Settings}
\label{tab:battery_settings}
\centering
\footnotesize
\setlength{\tabcolsep}{3.0pt}
\begin{tabular}{@{}>{\raggedright\arraybackslash}p{0.58\columnwidth}
                    >{\raggedright\arraybackslash}p{0.34\columnwidth}@{}}
\toprule
Parameter & Value \\
\midrule
Primary parameterization & Chen2020, LG INR21700-M50 \\
Additional parameterization & Ecker2015 \\
Chen2020 nominal capacity & 5.0 Ah \\
Chen2020 cell voltage range & 2.50--4.20 V \\
Chen2020 cell format & 21700 cylindrical, 70 $\times$ 21 mm \\
Initial/ambient temperature & 298.15 K \\
Maximum cell temperature & 328.15 K \\
\bottomrule
\end{tabular}
\end{table}

\begin{table}[!t]
\caption{Main Experimental Settings}
\label{tab:algorithm_settings}
\centering
\footnotesize
\setlength{\tabcolsep}{3.0pt}
\begin{tabular}{@{}>{\raggedright\arraybackslash}p{0.62\columnwidth}
                    >{\raggedright\arraybackslash}p{0.30\columnwidth}@{}}
\toprule
Setting & Value \\
\midrule
Total simulator budget & 56 \\
Initial evaluations & 6 \\
Number of matched trials & 5 \\
Initial LLM candidate pool & 15 \\
Regional-guidance horizon $T_{\mathrm L}$ & \RegionHorizon \\
Number of regional anchors $M$ & \RegionAnchors \\
Tchebycheff coefficient $\rho_{\mathrm{aug}}$ & \RhoAug \\
\bottomrule
\end{tabular}
\end{table}

ParEGO is the primary BO baseline because it uses the same decomposition--GP
structure as the numerical optimizer in LLMBO-MO~\cite{ref4}. DISK and PIMD
provide comparisons with surrogate-assisted multiobjective
evolutionary optimization~\cite{refDISK,refPIMD}; NSGA-II is included as a
classical evolutionary baseline~\cite{refNSGAII}.

\subsection{Performance Comparison}

LLMBO-MO obtains the highest final mean HV among the five methods on Chen2020
(Table~\ref{tab:chen_results}). After 56 evaluations, its mean HV is
\ChenLLMBOMean\,$\pm$\,\ChenLLMBOStd, compared with
\ChenParEGOMean\,$\pm$\,\ChenParEGOStd\ for ParEGO, a
\ChenFinalGain\ increase. NSGA-II, DISK, and PIMD reach mean HV values of
\ChenNSGAMean, \ChenDISKMean, and \ChenPIMDMean, respectively.
Figure~\ref{fig:chen_hv_5way} shows that the lead over ParEGO develops early
in the search and persists to the end of the budget.

\begin{figure}[!t]
\centering
\includegraphics[width=0.99\columnwidth]{chen2020_hv_5way.pdf}
\caption{HV trajectories of the five methods on Chen2020.}
\label{fig:chen_hv_5way}
\end{figure}

\begin{table}[!t]
\caption{Final HV on Chen2020}
\label{tab:chen_results}
\centering
\footnotesize
\setlength{\tabcolsep}{4.0pt}
\begin{tabular}{@{}lc@{}}
\toprule
Method & Final HV (mean $\pm$ SD) $\uparrow$ \\
\midrule
LLMBO-MO & \textbf{\ChenLLMBOMean\,$\pm$\,\ChenLLMBOStd} \\
ParEGO & \ChenParEGOMean\,$\pm$\,\ChenParEGOStd \\
NSGA-II & \ChenNSGAMean\,$\pm$\,\ChenNSGAStd \\
DISK & \ChenDISKMean\,$\pm$\,\ChenDISKStd \\
PIMD & \ChenPIMDMean\,$\pm$\,\ChenPIMDStd \\
\bottomrule
\end{tabular}
\end{table}

The comparison with ParEGO is repeated on Ecker2015. LLMBO-MO reaches
\EckerLLMBOMean\,$\pm$\,\EckerLLMBOStd\ at evaluation 56, whereas ParEGO
reaches \EckerParEGOMean\,$\pm$\,\EckerParEGOStd
(Table~\ref{tab:ecker_results}). Its HV is also higher over most of the search
in Fig.~\ref{fig:ecker_hv}. The advantage over ParEGO is therefore observed
with both battery parameterizations at the tested budget.

\begin{figure}[!t]
\centering
\includegraphics[width=0.99\columnwidth]{ecker2015_hv.pdf}
\caption{HV trajectories of LLMBO-MO and ParEGO on Ecker2015.}
\label{fig:ecker_hv}
\end{figure}

\begin{table}[!t]
\caption{HV Comparison on Ecker2015}
\label{tab:ecker_results}
\centering
\footnotesize
\setlength{\tabcolsep}{4.0pt}
\begin{tabular}{@{}lc@{}}
\toprule
Method & Evaluation 56 \\
\midrule
ParEGO & \EckerParEGOMean\,$\pm$\,\EckerParEGOStd \\
LLMBO-MO & \textbf{\EckerLLMBOMean\,$\pm$\,\EckerLLMBOStd} \\
\bottomrule
\end{tabular}
\end{table}

\subsection{Ablation Study}

The preceding comparisons assess the complete method. Two studies on
Chen2020 separate the effect of battery-specific prompting from the
contributions of initialization and regional guidance.

\subsubsection{LLM-Informed Initialization}

We compare random initialization, a detailed generic prompt, and a
battery-specific prompt (Figs.~\ref{fig:prompt_generic}
and~\ref{fig:prompt_warmstart}). All three variants use the same numerical BO
procedure after initialization. Each run has a budget of 26 evaluations,
and the results are computed over ten matched trials
(Table~\ref{tab:prompt_study}).

\begin{figure}[!t]
\centering
\promptbox{\footnotesize
You are an expert in lithium-ion battery fast-charging optimization. Task:
design $15$ diverse initial charging protocols for a multi-objective
Bayesian optimization workflow.\\
Decision variables and bounds: $I_1\in[2.0,6.0]$~A,
$I_2\in[2.0,5.0]$~A, $I_3\in[2.0,3.0]$~A,
$\Delta s_1\in[0.10,0.40]$, $\Delta s_2\in[0.10,0.30]$; hard constraint
$\Delta s_1+\Delta s_2<0.70$.\\
Objectives: minimize charging time, peak temperature rise, and capacity
fade.\\
Guidance: multi-objective optimization requires exploring different
trade-off directions (fast charging vs.\ thermal safety vs.\ aging
minimization); a good initial set should be well dispersed across the
five-dimensional decision space;
\textcolor{red!70!black}{consider leveraging physical relationships
between decision variables}; points near the boundary of the feasible
region can help maximize Pareto coverage, while interior points help model
the response surface in the central region. Respond with only a JSON
array.}
\caption{Generic prompt used in the initialization study of
Table~\ref{tab:prompt_study}. Unlike the battery-specific prompt in
Fig.~\ref{fig:prompt_warmstart}, it contains no battery-specific domain
knowledge.}
\label{fig:prompt_generic}
\end{figure}

\begin{table}[!t]
\caption{Effect of Different Initialization Strategies}
\label{tab:prompt_study}
\centering
\footnotesize
\setlength{\tabcolsep}{4.0pt}
\begin{tabular}{@{}lc@{}}
\toprule
Initialization & Final HV (mean $\pm$ SD) $\uparrow$ \\
\midrule
Random & 0.31062 $\pm$ 0.04303 \\
Generic prompt & 0.31371 $\pm$ 0.04675 \\
Battery-specific prompt & \textbf{0.36890 $\pm$ 0.00824} \\
\bottomrule
\end{tabular}
\end{table}

The generic prompt gives a mean final HV close to that of random
initialization: 0.31371 versus 0.31062. The battery-specific prompt reaches
0.36890, an increase of \PromptHVGain\ over random initialization, and has a
smaller standard deviation. The paired initialization analysis yields a
Holm-adjusted $p$-value of \PromptHolmP. The difference between generic and
battery-specific prompting indicates that the content of the domain
information matters for short-budget optimization, not merely the use of an
LLM to generate candidates.

\subsubsection{Regional Guidance}

Four variants are compared: the numerical BO optimizer without LLM guidance,
initialization only, regional guidance only, and the complete method. This
study uses three matched trials per variant. Each run contains six initial
evaluations followed by 50 BO iterations, giving a total budget of 56
battery-model evaluations. The results are summarized in
Table~\ref{tab:ablation} and Fig.~\ref{fig:ablation_hv_box}.

\begin{figure}[!t]

\centering

\includegraphics[width=0.99\columnwidth]{ablation_hv_box.pdf}

\caption{Final HV distributions of the four ablation variants on Chen2020.}

\label{fig:ablation_hv_box}

\end{figure}

\begin{table}[!t]

\caption{Ablation Results on Chen2020}

\label{tab:ablation}

\centering

\footnotesize

\setlength{\tabcolsep}{4.0pt}

\begin{tabular}{@{}lc@{}}

\toprule

Variant & Final HV (mean $\pm$ SD) $\uparrow$ \\

\midrule

No LLM & \AblationBaselineMean\,$\pm$\,\AblationBaselineStd \\

+ WarmStart & \AblationWarmMean\,$\pm$\,\AblationWarmStd \\

+ Regional guidance & \AblationRegionMean\,$\pm$\,\AblationRegionStd \\

LLMBO-MO & \textbf{\AblationFullMean\,$\pm$\,\AblationFullStd} \\

\bottomrule

\end{tabular}

\end{table}

Figure~\ref{fig:ablation_hv_box} and Table~\ref{tab:ablation} show that
LLM-informed initialization provides the largest improvement among the two
LLM-guidance components. Adding WarmStart increases the mean final HV from
\AblationBaselineMean\ to \AblationWarmMean, whereas regional guidance alone
reaches \AblationRegionMean. When regional guidance is added on top of
LLM-informed initialization, the mean final HV further increases from
\AblationWarmMean\ to \AblationFullMean. The complete method also has the
smallest sample standard deviation among the four variants.

The results indicate that the two components contribute at different stages
of the optimization. LLM-informed initialization improves the starting design
by directing the initial evaluations toward more promising charging protocols.
Region-Lift then uses the current scalarization weights to adjust candidate
selection during the subsequent BO iterations. In this comparison,
initialization provides most of the performance gain, while regional guidance
further improves the final optimization result.


\subsection{Comparison Across LLM Backends}

\label{sec:llm_backends}

To examine the effect of the LLM backend, Table~\ref{tab:llm_backends}
compares three models using the same Chen2020 prompt, numerical BO settings,
and evaluation budget. Each backend is evaluated in a single run consisting
of six initial evaluations followed by 50 BO iterations, for a total of 56
battery-model evaluations.

\begin{table}[!t]

\caption{HV Obtained with Different LLM Backends}

\label{tab:llm_backends}

\centering

\footnotesize

\setlength{\tabcolsep}{4.5pt}

\begin{tabular}{@{}lccc@{}}

\toprule

Backend & Initial HV & HV at 15 & Final HV \\

\midrule

DeepSeek-V4-Flash & 0.341225 & \textbf{0.373845} & 0.385986 \\

DeepSeek-V4-Pro & 0.345274 & 0.372526 & \textbf{0.387931} \\

GLM-5.3-Flash & \textbf{0.346248} & 0.372527 & 0.385324 \\

\bottomrule

\end{tabular}

\end{table}

The three backends produce different initial HV values, ranging from
0.341225 to 0.346248, indicating that the initial protocols generated by
the LLM depend to some extent on the backend. As the BO search proceeds,
however, the difference becomes smaller. At evaluation 15, the HV values
are concentrated between 0.372526 and 0.373845. After the full evaluation
budget, the final HV values are 0.385986, 0.387931, and 0.385324 for
DeepSeek-V4-Flash, DeepSeek-V4-Pro, and GLM-5.3-Flash, respectively.

Although the three LLMs produce somewhat different initial designs, their
optimization trajectories become closer after subsequent BO updates and
lead to similar final HV values. This suggests that the LLM backend mainly
affects the initial proposals, while the numerical BO procedure reduces
these differences as additional battery-model evaluations are collected.

\subsection{Analysis of Pareto-Optimal Charging Protocols}

HV evaluates the nondominated set as a whole. To examine the charging choices
represented by that set, we select protocols A--C from the Chen2020 Pareto
archive. Figure~\ref{fig:pareto_profiles} shows their objective values, and
Table~\ref{tab:representative_protocols} gives their currents and SOC
increments.

\begin{figure}[!t]
\centering
\includegraphics[width=\columnwidth]{pareto_protocols_ae.pdf}
\caption{Representative Pareto-optimal charging protocols on Chen2020.}
\label{fig:pareto_profiles}
\end{figure}

\begin{table*}[!t]
\caption{Representative Pareto-Optimal Charging Protocols}
\label{tab:representative_protocols}
\centering
\footnotesize
\setlength{\tabcolsep}{5.0pt}
\begin{tabular}{@{}c l r r r@{}}
\toprule
Protocol & $I_1/I_2/I_3/\Delta s_1/\Delta s_2/\Delta s_3$ &
$t_c$ (s) & $T_p$ (K) & $Q_s$ \\
\midrule
A & 6.000/5.000/3.000/0.400/0.300/0.100 & 2880 & 7.567 & 1.261 \\
B & 2.782/2.000/3.000/0.211/0.100/0.489 & 5199 & 3.069 & 0.680 \\
C & 2.000/2.000/2.000/0.245/0.153/0.403 & 7200 & 1.528 & 0.640 \\
\bottomrule
\end{tabular}
\end{table*}

Protocol A applies higher currents early in charging and reaches the target
SOC in 2880~s. It is the fastest of the three, but has the largest peak
temperature rise and aging indicator. Protocol C uses a constant 2~A current
and takes 7200~s, with the lowest values of both indicators. Protocol B offers
an intermediate choice: relative to C, it reduces charging time to 5199~s,
while the aging indicator increases from 0.640 to 0.680 and peak temperature
rise from 1.528 to 3.069~K. The benefit of faster charging can thus be weighed
against its thermal and degradation costs using the returned protocols.

Figure~\ref{fig:abc_charging_profiles} shows the voltage, temperature,
current, and SOC trajectories. The current transitions distinguish the
multistage protocols A and B from the constant-current protocol C, and the SOC
curves show their different completion times. The voltage trajectories remain
below the prescribed limit.

\begin{figure}[!t]
\centering
\includegraphics[width=\columnwidth]{abc_charging_profiles.pdf}
\caption{Charging trajectories of protocols A--C on Chen2020.}
\label{fig:abc_charging_profiles}
\end{figure}
